% Vietnamese translation for OI Public Library
% Source-PDF: IOI中国国家候选队论文/2007/day1/5.杨沐《浅析信息学中的“分”与“合”》.pdf
\documentclass[11pt]{article}
\usepackage{oipl-vn}

\newcommand{\Oh}{\mathcal{O}}
\newcommand{\dist}{\operatorname{dist}}
\newcommand{\deep}{\operatorname{deep}}
\newcommand{\len}{\operatorname{len}}

\title{Bàn về tư duy ``phân'' và ``hợp'' trong tin học}
\author{Yang Mu\\Trường Trung học số 3 Phúc Châu, Phúc Kiến}
\date{2007}

\begin{document}
\maketitle

\origtitle{An elementary analysis of ``division'' and ``combination'' in informatics}
\sourcepath{IOI China candidate team papers / 2007 / day 1 / Yang Mu}

\section*{Tóm tắt}

Bài viết thảo luận việc vận dụng tư duy ``phân'' và ``hợp'' trong tin học. Cốt
lõi của tư duy này là tìm ra bản chất của bài toán qua việc chuyển đổi qua lại
giữa tách nhỏ và gom lại, từ đó giải quyết bài toán. Phạm vi của nó không chỉ
giới hạn trong ``chia để trị'' theo nghĩa truyền thống. Thông qua ba ví dụ
\textit{Mẫu sữa}, \textit{Tái dựng cây} và \textit{Dãy tối ưu}, bài viết giới
thiệu cách dùng cụ thể của tư duy ``phân'' và ``hợp'' trong quá trình giải
bài.

\paragraph{Từ khóa.}
Tư duy phân và hợp; chuyển hóa; quan hệ biện chứng.

\section{Dẫn nhập}

``Thế lớn trong thiên hạ, hợp lâu tất phân, phân lâu tất hợp'' là câu mở đầu
giàu ý vị biện chứng và có thể xem là tinh thần của \textit{Tam quốc diễn
nghĩa}. Lịch sử lâu dài cũng nhiều lần cho thấy sự chuyển hóa giữa hợp và
phân: một quá trình phát triển xoáy ốc, từ thấp đến cao, trong đó hai mặt ấy
liên tục tác động lẫn nhau.

Quy luật này không chỉ xuất hiện trong đời sống xã hội. Trong nhiều bài toán
tin học, cách đi từ trạng thái ``phân'' sang ``hợp'', hoặc ngược lại, cũng bộc
lộ rất rõ trong quá trình tìm lời giải.

Tư duy ``phân'' là chia một bài toán lớn, khó xử lý trực tiếp, thành những bài
toán con nhỏ hơn hoặc thành các trường hợp có thêm ràng buộc, rồi từ đó tìm
đường giải. Tư duy ``hợp'' đi theo chiều ngược lại: gom các mảnh nhỏ đã tách
ra thành một bài toán lớn hơn, qua đó làm rõ quan hệ giữa chúng và rút ra
phương pháp.

Trong các kỳ thi tin học, ta thường gặp bài toán có quy mô quá lớn nên khó
quan sát, hoặc có quá nhiều tình huống khiến việc suy nghĩ trở nên nặng nề.
Khi ấy, thử nhìn bài toán bằng cặp tư duy ``phân'' và ``hợp'' thường giúp tìm
được điểm vào.

\section{Phân tích ví dụ}

Sau đây là ba ví dụ: \textit{Mẫu sữa}, \textit{Tái dựng cây} và \textit{Dãy
tối ưu}. Mỗi ví dụ dùng một khía cạnh khác nhau của tư duy ``phân'' và
``hợp''.

\subsection{Ví dụ 1: Mẫu sữa}

\paragraph{Phát biểu.}
Do chất lượng sữa của đàn bò dao động từng ngày, Farmer John muốn tìm quy
luật dao động ấy. Ông ghi lại một dãy
\[
  a_1,a_2,\ldots,a_N
\]
gồm các số nguyên từ \(0\) đến \(1\,000\,000\), với
\(1\le N\le 20\,000\). Cần tìm độ dài lớn nhất của một đoạn mẫu xuất hiện
trong dãy ít nhất \(k\) lần. Ví dụ, trong dãy \(1,2,3,2,3,2,3,1\), đoạn
\(2,3,2,3\) xuất hiện hai lần, nên độ dài mẫu lớn nhất là \(4\).

Bài này lấy từ USACO December 2006 Gold, \textit{Milk Patterns}.

\paragraph{Phân tích ban đầu.}
Đây là bài toán liên quan đến so khớp chuỗi. Một công cụ tự nhiên là cây hậu
tố; thật vậy, cây hậu tố giải được bài toán trong thời gian tuyến tính. Tuy
nhiên độ phức tạp cài đặt cao khiến nó không phải lựa chọn thuận tiện trong
thời gian thi. Ta cần tìm một cách đơn giản hơn mà vẫn đủ hiệu quả.

\paragraph{Liệt kê trực tiếp.}
Cách đơn giản nhất là xét mọi đoạn con của dãy, so sánh chúng với nhau, đếm
số lần xuất hiện của từng đoạn rồi lấy đoạn dài nhất xuất hiện ít nhất \(k\)
lần. Số đoạn con đạt cỡ \(N^2\), còn việc so sánh hai đoạn có thể tốn
\(\Oh(N)\), vì vậy cách này có độ phức tạp đến \(\Oh(N^5)\) nếu làm thô theo
từng cặp và từng vị trí. Nó hoàn toàn không đáp ứng được giới hạn.

Nguyên nhân chính của sự kém hiệu quả là có quá nhiều phép so sánh lặp lại.
Ta có thể nghĩ đến KMP: sau khi chọn một đoạn con làm mẫu, dùng KMP để đếm số
lần nó xuất hiện trong dãy trong \(\Oh(N)\) thời gian. Khi đó tổng độ phức tạp
còn \(\Oh(N^3)\). Dù vậy, việc tiếp tục cải thiện bằng các thuật toán so khớp
chuỗi trực tiếp gặp nhiều trở ngại. Ta quay lại điểm xuất phát và đổi góc nhìn.

\paragraph{Cải tiến cách xét.}
Quan sát cách làm thô: việc xét mọi đoạn con có nhiều phần dư thừa, bởi hai
đoạn chỉ có thể bằng nhau nếu chúng cùng độ dài. Gợi ý này dẫn đến việc xét
theo độ dài đoạn.

Bài toán có một tính đơn điệu rõ ràng. Nếu tồn tại đoạn con \(S_{\len}\) độ
dài \(\len\) xuất hiện ít nhất \(k\) lần, thì với mọi
\(\len'<\len\), cũng tồn tại một đoạn con \(S_{\len'}\) độ dài \(\len'\) xuất
hiện ít nhất \(k\) lần: chỉ cần lấy cùng một vị trí tương đối trong mỗi bản
sao của \(S_{\len}\). Vì vậy ta có thể tìm kiếm nhị phân đáp án. Với một độ
dài cố định \(\len\), chỉ cần kiểm tra xem có đoạn độ dài \(\len\) nào xuất
hiện ít nhất \(k\) lần hay không.

Sau khi đổi cách xét, nút thắt mới là kiểm tra hai đoạn cùng độ dài có bằng
nhau không. Với độ dài \(\len\), chỉ có \(N-\len+1\) đoạn ứng viên. Ta lưu
chúng trong bảng băm và dùng giá trị băm để nhận diện.

Với đoạn
\[
  \mathrm{sub}_t=(a_t,a_{t+1},\ldots,a_{t+\len-1}),
\]
chọn hai số nguyên tố \(p_1,p_2\) và hai mô-đun lớn
\(\mathrm{hashsize}_1,\mathrm{hashsize}_2\). Đặt
\[
\begin{aligned}
  f_t &=
  \left(\sum_{j=0}^{\len-1} a_{t+j}p_1^{\len-1-j}\right)
  \bmod \mathrm{hashsize}_1,\\
  \mathrm{pos}_t &=
  \left(\sum_{j=0}^{\len-1} a_{t+j}p_2^{\len-1-j}\right)
  \bmod \mathrm{hashsize}_2.
\end{aligned}
\]
Giá trị \(\mathrm{pos}_t\) cho biết vị trí trong bảng băm, còn \(f_t\) thay
cho nội dung đoạn khi cần xử lý va chạm. Nếu hai đoạn có cả hai giá trị
\(\mathrm{pos}\) và \(f\) giống nhau, ta xem chúng là cùng một đoạn. Với hai
hàm băm độc lập, xác suất nhầm lẫn rất nhỏ trong bối cảnh bài này.

Do hai đoạn liên tiếp \(\mathrm{sub}_t\) và \(\mathrm{sub}_{t+1}\) chỉ khác
nhau ở một phần tử bị bỏ ở đầu và một phần tử thêm ở cuối, các giá trị băm
có thể cập nhật cuốn chiếu:
\[
  f_{t+1}=(f_t p_1-a_t p_1^{\len}+a_{t+\len})
  \bmod \mathrm{hashsize}_1,
\]
và \(\mathrm{pos}_{t+1}\) được tính tương tự với \(p_2\). Vì vậy mọi đoạn độ
dài \(\len\) được xử lý trong \(\Oh(N)\) thời gian. Kết hợp với tìm kiếm nhị
phân, độ phức tạp kỳ vọng là \(\Oh(N\log N)\).

\paragraph{Khung thuật toán.}
\begin{verbatim}
check(len):
    clear hash table
    compute rolling hashes for all length-len substrings
    insert each pair (pos, f) and count its occurrences
    return true if some count reaches k

answer:
    binary search len from 0 to N using check(len)
\end{verbatim}

\paragraph{Nhận xét.}
Trong ví dụ này, tư duy ``phân'' xuất hiện qua việc tách bài toán tối ưu dưới
các ràng buộc \(N,k\) thành bài toán tồn tại dưới các ràng buộc \(N,k,\len\).
Việc thêm tham số \(\len\) làm cấu trúc bài toán rõ hơn, từ đó cho ra một
thuật toán ngắn gọn và hiệu quả.

\subsection{Ví dụ 2: Tái dựng cây}

\paragraph{Phát biểu.}
Cho số lá \(l\) của một cây, với \(2\le l\le 200\), và khoảng cách trong cây
giữa mọi cặp lá. Mọi cạnh của cây có trọng số \(1\), mọi khoảng cách không
vượt quá \(200\). Hãy tái dựng một cây phù hợp. Các đỉnh chưa biết có thể được
đặt nhãn tùy ý. Nếu không tồn tại cây như vậy, in \(-1\).

Bài này lấy từ ZOJ 2701, \textit{Restore the Tree}.

\paragraph{Phân tích các trường hợp nhỏ.}
Điều kiện đầu vào chỉ mô tả quan hệ khoảng cách giữa các lá; các đỉnh trong
không được mô tả trực tiếp, nên dựng cây ngay từ đầu rất khó. Ta bắt đầu bằng
các trường hợp nhỏ.

Nếu chỉ có hai lá, đáp án hiển nhiên là một đường đi đơn giữa hai lá ấy:
\[
  1 \quad \text{-- -- -- -- --} \quad 2.
\]

Với ba lá, đặt \(\dist[i,j]\) là khoảng cách giữa lá \(i\) và lá \(j\). Để
thuận tiện, lấy lá \(1\) làm gốc suy nghĩ.

Trước hết, nếu
\[
  \dist[2,3]>\dist[1,2]+\dist[1,3],
\]
thì vi phạm bất đẳng thức tam giác trên cây, nên vô nghiệm. Hai bất đẳng thức
tương tự với các cặp còn lại cũng cho vô nghiệm.

Nếu
\[
  \dist[2,3]=\dist[1,2]+\dist[1,3],
\]
thì lá \(1\) nằm trên đường đi từ \(2\) đến \(3\). Khi đó \(1\) không thể là lá
trong cây cần dựng, mâu thuẫn với đề bài. Các trường hợp bằng nhau tương tự
cũng vô nghiệm.

Trường hợp còn lại là các bất đẳng thức đều nghiêm. Khi ấy ba đường đi từ ba
lá gặp nhau tại một điểm phân nhánh. Độ sâu của điểm phân nhánh tính từ lá
\(1\) là
\[
  x=\frac{\dist[1,2]+\dist[1,3]-\dist[2,3]}{2}.
\]
Do đó nếu \(\dist[1,2]+\dist[1,3]-\dist[2,3]\) là số lẻ, \(x\) không nguyên
và bài toán vô nghiệm. Nếu chẵn, lá \(3\) nằm ở cuối một nhánh tách ra từ
đỉnh thứ \(x\) trên đường đi từ lá \(1\) đến lá \(2\), với độ dài nhánh là
\(\dist[1,3]-x\).

Có thể hình dung dạng cây ba lá như sau:
\[
\begin{array}{ccccc}
&&1&&\\
&&|&&\\
&&\vdots&&\\
&&s&&\\
&\swarrow&&\searrow&\\
2&&&&3
\end{array}
\qquad
x=\deep(s).
\]

\paragraph{Phân tích tiếp.}
Từ trường hợp ba lá, ta rút ra một nhận xét quan trọng: hai lá bất kỳ luôn
được nối bởi một đường đi gồm các đỉnh trong, và vị trí của mọi lá còn lại
so với đường đi ấy đều có thể xác định bằng công thức ba lá ở trên.

Gọi bài toán với \(n\) lá là \(P(n)\). Nếu biết cách ghép một lá mới vào một
cây đã dựng, thì \(P(n)\) chuyển thành \(P(n-1)\). Cụ thể, giả sử cây hiện tại
\(T\) đã chứa \(k-1\) lá, mang nhãn \(1,2,\ldots,k-1\). Lấy lá \(1\) làm gốc,
và cần ghép thêm lá \(k\).

Với mỗi lá \(i\ne 1\) đã có trong \(T\), xét ba lá \(1,i,k\). Công thức trên
cho ta điểm phân nhánh \(s_i\) của lá \(k\) đối với đường đi \(1\to i\), cùng
độ sâu \(\deep(s_i)\):
\[
  \deep(s_i)=
  \frac{\dist[1,i]+\dist[1,k]-\dist[i,k]}{2}.
\]
Nếu cây thật tồn tại, mọi điểm \(s_i\) này đều nằm trên cùng đường đi từ \(1\)
đến \(k\). Khi ghép lá mới, phần cây mới duy nhất là đoạn từ điểm phân nhánh
sâu nhất đến lá \(k\). Đặt
\[
  s_{\max}=\arg\max_i \deep(s_i).
\]
Ta tách từ \(s_{\max}\) một đường mới dài
\[
  \dist[1,k]-\deep(s_{\max})
\]
và đặt cuối đường ấy là lá \(k\).

Quan sát này cũng giúp giảm mạnh gánh nặng xét trường hợp. Trong quá trình
ghép, số tình huống vô nghiệm khá nhiều; nếu xử lý hết ngay tại từng bước,
tư duy và cài đặt đều phức tạp. Nhưng nếu đề có nghiệm, cách ghép trên sẽ đưa
lá mới vào đúng vị trí. Vì vậy ta có thể dựng một cây ứng viên trước, rồi
kiểm tra toàn bộ ở cuối.

\paragraph{Khung thuật toán.}
\begin{enumerate}
  \item Dựng cây cơ sở \(T\), ban đầu chỉ gồm lá \(1\), lá \(2\), và một đường
        đi gồm \(\dist[1,2]\) cạnh nối chúng.
  \item Lấy một lá chưa có trong \(T\). Với mỗi lá cũ \(i\ne 1\), tính điểm
        phân nhánh \(s_i\) theo công thức ba lá. Chọn điểm sâu nhất
        \(s_{\max}\), rồi gắn lá mới bằng một đường có độ dài
        \(\dist[1,k]-\deep(s_{\max})\).
  \item Lặp đến khi mọi lá đã được gắn.
  \item Kiểm tra cây ứng viên: số lá phải đúng bằng \(l\), và với mọi cặp lá
        \((i,j)\), khoảng cách trong cây phải bằng dữ liệu \(\dist[i,j]\).
        Nếu kiểm tra thất bại thì vô nghiệm; ngược lại cây thu được là một
        nghiệm hợp lệ.
\end{enumerate}

Với \(l\le 200\) và khoảng cách tối đa \(200\), số đỉnh tạo ra vẫn cùng bậc với
tổng độ dài các nhánh. Cách cài đặt thông thường duy trì độ sâu, cha và bảng
vị trí trên các đường đi cho độ phức tạp thời gian và bộ nhớ \(\Oh(l^2)\),
đạt đến mức cần thiết vì bản thân dữ liệu khoảng cách đã có kích thước
\(\Theta(l^2)\).

\paragraph{Nhận xét.}
Ở bài này, ta đi từ trường hợp nhỏ để phát hiện tính chất bản chất, rồi dùng
``phân'' và ``hợp'' để chuyển bài toán \(n\) lá thành bài toán \(n-1\) lá cộng
với một thao tác ghép lá. Chính việc gom lại ở cuối bằng bước kiểm tra toàn cục
cho phép bỏ qua nhiều nhánh xét riêng lẻ trong quá trình dựng.

\subsection{Ví dụ 3: Dãy tối ưu}

\paragraph{Phát biểu.}
Cho một dãy số nguyên dương độ dài \(n\), với \(n\le 1000\). Cần chọn một dãy
con sao cho trong mọi đoạn liên tiếp độ dài \(m\), số phần tử được chọn không
vượt quá \(k\), trong đó \(k,m\le 100\). Tổng các phần tử được chọn phải lớn
nhất.

\paragraph{Phân tích ban đầu.}
Thoạt nhìn, quy hoạch động là công cụ tự nhiên cho bài toán tối ưu kiểu này.
Nhưng ràng buộc trên mọi đoạn độ dài \(m\) làm trạng thái trở nên nặng. Nếu
nén trạng thái theo \(m\) vị trí gần nhất thì độ phức tạp có thể lên tới
\(\Oh(2^m n)\), không thể dùng. Cây đoạn cũng phù hợp với thao tác trên đoạn
một chiều, nhưng sau nhiều thử nghiệm vẫn không xử lý được ràng buộc toàn cục
của bài toán.

Khó khăn chính nằm ở việc \(k\) và \(m\) đều khá lớn, khiến quan hệ giữa các
lựa chọn quá phức tạp. Ta tách bài toán thành một phiên bản đơn giản hơn.

\paragraph{Tách phức tạp thành đơn giản.}
Xét bài toán con: chọn một dãy con sao cho trong mọi đoạn độ dài \(m\), có
không quá \(1\) phần tử được chọn, và tổng là lớn nhất.

Bài toán con này có quy hoạch động rất gọn. Đặt \(dp[i]\) là tổng lớn nhất có
thể lấy trong đoạn \([1,i]\). Khi xét phần tử \(i\), hoặc không lấy nó, hoặc
lấy nó và khi đó phần tử được chọn trước đó phải nằm trước \(i-m+1\):
\[
\begin{aligned}
  dp[i]&=\max(dp[i-1], dp[i-m]+\mathrm{value}[i]) && (i\ge m),\\
  dp[i]&=\max(dp[i-1], \mathrm{value}[i]) && (0<i<m),\\
  dp[0]&=0.
\end{aligned}
\]

Câu hỏi là bài toán con này liên hệ thế nào với bài toán gốc.

\paragraph{Một mệnh đề cấu trúc.}
Một nghiệm của bài toán gốc tương đương với hợp của \(k\) nghiệm bài toán con
đôi một không giao nhau.

\paragraph{Định lý 1.}
Mọi nghiệm của bài toán gốc đều có thể được phân hoạch thành \(k\) nghiệm con
không giao nhau của phiên bản ``mỗi đoạn độ dài \(m\) chọn không quá \(1\)
phần tử''.

\paragraph{Chứng minh.}
Lấy một nghiệm bất kỳ
\[
  P=\{a_1,a_2,\ldots,a_t\},\qquad
  a_1<a_2<\cdots<a_t.
\]
Gọi \(\mathrm{sum}[i]\) là số phần tử được chọn trong đoạn \([1,i]\). Theo đề
bài,
\[
  \mathrm{sum}[i]-\mathrm{sum}[i-m]\le k \qquad (i\ge m).
\]
Từ đó suy ra với mọi \(j>k\),
\[
  a_j-a_{j-k}\ge m.
\]
Bây giờ chia \(P\) thành \(k\) nhóm theo chỉ số modulo \(k\):
\[
  P_r=\{a_j: j\equiv r \pmod k\},\qquad 0\le r<k.
\]
Trong mỗi nhóm \(P_r\), hai phần tử liên tiếp cách nhau ít nhất \(m\), nên mỗi
đoạn độ dài \(m\) chứa không quá một phần tử của nhóm ấy. Vì vậy mỗi \(P_r\)
là nghiệm hợp lệ của bài toán con, và các nhóm không giao nhau, hợp lại đúng
bằng \(P\).

\paragraph{Định lý 2.}
Hợp của \(k\) nghiệm con đôi một không giao nhau luôn là một nghiệm hợp lệ của
bài toán gốc.

\paragraph{Chứng minh.}
Gọi các nghiệm con là \(P_0,P_1,\ldots,P_{k-1}\). Trong mỗi \(P_i\), hai phần
tử được chọn liên tiếp cách nhau ít nhất \(m\), nên một đoạn bất kỳ độ dài
\(m\) chứa nhiều nhất một phần tử của \(P_i\). Vì thế đoạn ấy chứa nhiều nhất
\(k\) phần tử của
\[
  P=P_0\cup P_1\cup\cdots\cup P_{k-1}.
\]
Do đó \(P\) thỏa ràng buộc của bài toán gốc.

Hai định lý trên cho thấy điều kiện cần và đủ: nghiệm gốc chính là cách gom
\(k\) nghiệm con không giao nhau.

\paragraph{Từ bài toán con đến mô hình tổng thể.}
Ta đã chuyển ``một nghiệm tối ưu'' thành ``hợp tối ưu của \(k\) nghiệm con''.
Vấn đề còn lại là tìm được \(k\) nghiệm con không giao nhau có tổng lớn nhất.

Không thể chỉ chạy quy hoạch động bài toán con \(k\) lần, vì các nghiệm con
có thể chọn trùng một phần tử. Tham lam theo từng nghiệm con cũng không đúng:
\(k\) nghiệm tối ưu cục bộ khi hợp lại chưa chắc tạo thành nghiệm toàn cục.

Đây là lúc cần tư duy ``hợp''. Thay vì xử lý từng nghiệm con riêng biệt, ta
xây dựng một mô hình dòng chi phí cực đại để xét đồng thời mọi nghiệm con.
Những bài toán có ràng buộc ``mỗi điểm được chọn tối đa một lần'' thường phù
hợp với mạng dòng.

\paragraph{Mô hình dòng.}
Dựng mạng có hướng có trọng số \(D=(V,A,C)\):
\begin{itemize}
  \item Mỗi phần tử \(i\) của dãy được biểu diễn bằng hai đỉnh \(i\) và \(i'\).
        Thêm cung \(i\to i'\) có dung lượng \(1\), chi phí
        \(\mathrm{value}[i]\). Đi qua cung này tương ứng với chọn phần tử \(i\).
  \item Với mọi \(i\), thêm cung \(i\to i+1\) có dung lượng vô hạn và chi phí
        \(0\). Cung này cho phép bỏ qua phần tử.
  \item Nguồn \(S\) nối đến mọi đỉnh \(i\) bằng cung dung lượng vô hạn, chi phí
        \(0\).
  \item Mọi đỉnh \(i'\) nối đến đích \(T\) bằng cung dung lượng vô hạn, chi
        phí \(0\).
  \item Với mọi \(i\) sao cho \(i+m\le n\), thêm cung \(i'\to i+m\) có dung
        lượng vô hạn, chi phí \(0\). Cung này bảo đảm sau khi chọn \(i\), cùng
        một đơn vị dòng chỉ có thể chọn tiếp từ vị trí cách ít nhất \(m\).
\end{itemize}

Sơ đồ khái niệm là
\[
\begin{array}{cccccc}
&&S&&&\\
&\swarrow&\downarrow&\searrow&&\\
1&\to&2&\to&\cdots&\to n\\
\downarrow&&\downarrow&&&\downarrow\\
1'&\to&2'&\to&\cdots&\to n'\\
&\searrow&\downarrow&\swarrow&&\\
&&T&&&
\end{array}
\]
trong đó các cung \(i\to i'\) mang lợi ích \(\mathrm{value}[i]\), còn cung
\(i'\to i+m\) thể hiện khoảng cách tối thiểu giữa hai phần tử do cùng một
nghiệm con chọn.

Mỗi đơn vị dòng từ \(S\) đến \(T\) tương ứng với một nghiệm con của phiên bản
``không quá một phần tử trong mọi đoạn độ dài \(m\)''. Vì cung \(i\to i'\) có
dung lượng \(1\), một phần tử không thể được chọn bởi hai nghiệm con khác
nhau. Do đó chỉ cần tìm \(k\) lần đường tăng chi phí lớn nhất, hoặc tương
đương gửi \(k\) đơn vị dòng có tổng chi phí lớn nhất, ta thu được đáp án.

Đồ thị có số cung cùng bậc với số đỉnh. Nếu dùng SPFA để tìm đường tăng theo
chi phí lớn nhất trong mạng dư, độ phức tạp kỳ vọng khoảng \(\Oh(kN)\), rất
tốt với giới hạn của bài.

\paragraph{Khung thuật toán.}
\begin{verbatim}
build graph:
    for i = 1..n:
        add edge i -> i' with capacity 1 and cost value[i]
        add edge S -> i with infinite capacity and cost 0
        add edge i' -> T with infinite capacity and cost 0
        if i < n: add edge i -> i+1 with infinite capacity and cost 0
        if i+m <= n: add edge i' -> i+m with infinite capacity and cost 0

answer:
    repeat k times:
        find a maximum-cost augmenting path from S to T
        augment one unit of flow
    return total cost
\end{verbatim}

\paragraph{Nhận xét.}
Khi bài toán khó tiếp cận, ta dùng ``phân'' để đơn giản hóa, qua đó nhận ra
cấu trúc của nghiệm. Sau đó dùng ``hợp'' để đưa mọi bài toán con vào một mô
hình dòng duy nhất. Ví dụ này cho thấy vai trò rất mạnh của hai thao tác ấy
trong tư duy giải bài.

\section{Kết luận}

Bài viết đã giới thiệu vài cách dùng tư duy ``phân'' và ``hợp'': khi bài toán
có tính đơn điệu, tách theo tham số để dùng tìm kiếm nhị phân; khi phát hiện
tính chất quy nạp, bắt đầu từ phạm vi nhỏ rồi mở rộng dần; và khi bài toán
khó quan sát trực tiếp, đơn giản hóa để tìm cấu trúc rồi gom các phần lại
bằng một mô hình tổng thể.

Trong ví dụ \textit{Mẫu sữa}, còn nhiều lời giải hay khác, nhưng cách tiếp cận
từ ``phân'' có độ khó tư duy thấp, dễ cài đặt và phù hợp với môi trường thi
cần tốc độ. Trong ví dụ \textit{Tái dựng cây}, quy về bài toán nhỏ hơn gần như
là con đường bắt buộc. Hai ví dụ đầu cũng cho thấy ``phân'' và ``hợp'' tuy đối
lập nhưng không có ranh giới cứng: đã phân thì phải có lúc hợp lại, và trong
hợp vẫn có phân.

Tư duy ``phân'' giúp nhanh chóng đi vào lõi bài toán, nhưng nếu chia quá vụn
thì vấn đề trở nên rời rạc và mất phương hướng. Tư duy ``hợp'' xâu các phần
phức tạp lại thành một chỉnh thể, làm nổi bật quan hệ giữa chúng. Hai mặt ấy
bổ sung cho nhau.

Những ví dụ trong bài chỉ chạm đến một phần nhỏ của chủ đề. Chia để trị,
chuyển hóa qua phần bù và nhiều cấu trúc dữ liệu cao cấp đều là những ứng
dụng quen thuộc của cùng một tinh thần. Quan trọng hơn, ``phân'' và ``hợp''
không chỉ là kỹ thuật cài đặt mà còn là phương pháp tư duy. Dù tách hay gom,
mục tiêu đều là chuyển bài toán về hình thức dễ suy nghĩ hơn. Kết hợp kiến
thức sẵn có với tư duy chuyển hóa, đồng thời tích lũy kinh nghiệm và dám thử
cách nhìn mới, ta mới có thể tìm ra bản chất trong các bài toán biến hóa và
giải chúng nhanh hơn, hiệu quả hơn.

\section*{Lời cảm ơn}

Tác giả cảm ơn cô Wei Lizhen đã hướng dẫn và giúp đỡ trong quá trình viết bài;
cảm ơn anh Wang Zhikun đã đọc kỹ và góp ý chỉnh sửa; cảm ơn bạn Yuan Xinhao
đã cung cấp ví dụ thứ ba; và cảm ơn bạn Wang Hang vì sự giúp đỡ trong quá
trình viết.

\section*{Tài liệu tham khảo}

\begin{enumerate}[label={[\Alph*]}]
  \item Liu Rujia, Huang Liang, \textit{Nghệ thuật thuật toán và thi tin học}.
  \item Yu Yuanming, \textit{Thuật toán đường đi ngắn nhất và ứng dụng}, luận
        văn trại đông 2006.
\end{enumerate}

\appendix
\section{Đề bài Mẫu sữa}

Farmer John nhận thấy chất lượng sữa của đàn bò thay đổi từng ngày. Sau khi
tìm hiểu thêm, ông phát hiện tuy không thể dự đoán chính xác chất lượng sữa
của ngày kế tiếp, dữ liệu vẫn có những mẫu lặp đều đặn.

Mỗi mẫu sữa được ghi thành một số nguyên từ \(0\) đến \(1\,000\,000\). Với dữ
liệu của một con bò trong \(N\) ngày, \(1\le N\le 20\,000\), hãy tìm mẫu dài
nhất xuất hiện giống hệt ít nhất \(K\) lần, với \(2\le K\le N\). Các lần xuất
hiện có thể chồng lấn; chẳng hạn dãy \(1,2,3,2,3,2,3,1\) chứa mẫu
\(2,3,2,3\) hai lần. Dữ liệu bảo đảm có ít nhất một mẫu xuất hiện ít nhất
\(K\) lần.

\paragraph{Dữ liệu vào.}
Dòng đầu chứa \(N\) và \(K\). \(N\) dòng tiếp theo chứa chất lượng sữa từng
ngày.

\paragraph{Dữ liệu ra.}
In một số nguyên: độ dài mẫu dài nhất xuất hiện ít nhất \(K\) lần.

\paragraph{Ví dụ.}
\[
\begin{array}{ll}
\text{Input:}&
\begin{array}[t]{l}
8\ 2\\
1\\2\\3\\2\\3\\2\\3\\1
\end{array}\\[1em]
\text{Output:}&4
\end{array}
\]

\section{Đề bài Tái dựng cây}

Một cây là đồ thị vô hướng liên thông không có chu trình. Đỉnh có bậc \(1\)
trong cây được gọi là lá. Với một cây, ta có thể xác định khoảng cách giữa
mọi cặp lá, tức số cạnh trên đường đi giữa hai lá ấy. Cho tất cả các khoảng
cách này, hãy tái dựng cây.

Dữ liệu gồm nhiều bộ. Mỗi bộ bắt đầu bằng \(l\), số lá của cây, với
\(2\le l\le 200\). Tiếp theo là ma trận \(l\times l\) khoảng cách giữa các lá.
Mọi khoảng cách không âm, không vượt quá \(200\), đường chéo chính bằng \(0\),
và ma trận đối xứng.

Nếu không tồn tại cây phù hợp, in \(-1\). Ngược lại, in số đỉnh \(n\), rồi in
\(n-1\) cạnh của cây. \(l\) đỉnh đầu tiên phải tương ứng với các lá trong dữ
liệu; các đỉnh khác được đánh số tùy ý. Nếu có nhiều nghiệm, in nghiệm bất kỳ.

\paragraph{Ví dụ.}
\[
\begin{array}{ll}
\text{Input:}&
\begin{array}[t]{l}
3\\
0\ 2\ 3\\
2\ 0\ 3\\
3\ 3\ 0
\end{array}\\[1em]
\text{Output:}&
\begin{array}[t]{l}
5\\
1\ 4\\
2\ 4\\
3\ 5\\
4\ 5
\end{array}
\end{array}
\]

\end{document}
